\section{Conclusions}
\label{sec:conclusions}

We demonstrated a representation-level diagnostic of observed--injected
differences in selected DESI Legacy Imaging Surveys DR10 images. Linear probes
reuse frozen trained features to evaluate this separation directly.
Strong held-out discrimination in the operational finder's penultimate
representation is accompanied by discrimination in survey self-supervised
features, generic ImageNet features, and elementary pixel statistics under the
common dependency-aware procedure. The separation appears in the operational
detector, two external learned representations, and simple pixel statistics.

A leading pixel-statistic offset largely predates insertion, and the
separately fitted operational observed--bare axis has a near-zero mean signed
movement with heterogeneous individual shifts.

Separately, the finder recovers 3.45 per cent of the full grid without a
visibility cut.
Modelled source-photon sampling lowers this response, and the tested paired
single-S\'{e}rsic construction has higher recovery than its regularized
structured counterpart. The brightness control shows how the same targets
respond to source scaling, while the checkpoint panel shows which intervention
results persist across trained models. The combined results show why future
injection pipelines should be compared on matched foreground populations and
observing conditions, with bare-target controls used to separate pre-existing
image differences from the effects of source insertion.
